\documentclass[10pt,a4paper]{amsart}
\usepackage{amsmath,amssymb,amsthm,graphicx}
% Retained preamble for the Stacks Project chapter snapshot in this directory.
% The chapter's own preamble.tex is a copy of the shared Stacks Project preamble; the block
% below is sliced MECHANICALLY, line for line, from it. Nothing was added, removed or
% redefined; the wrapper class, the xr-hyper/hyperref wiring and the externaldocument list
% are omitted because the assembler supplies the document class.
%
% lines taken from preamble.tex: 3, 5, 9, 11, 12, 156, 157, 158, ... (51 lines total)
%
% package coverage restored from the chapter's own preamble (2026-10-13): xy, xy-2cell

\usepackage[all]{xy}
\xyoption{2cell}
\UseAllTwocells

\newenvironment{reference}{\comment}{\endcomment}
\newenvironment{slogan}{\comment}{\endcomment}
\newenvironment{history}{\comment}{\endcomment}
\theoremstyle{plain}
\newtheorem{theorem}[subsection]{Theorem}
\newtheorem{proposition}[subsection]{Proposition}
\newtheorem{lemma}[subsection]{Lemma}
\theoremstyle{definition}
\newtheorem{definition}[subsection]{Definition}
\newtheorem{example}[subsection]{Example}
\newtheorem{exercise}[subsection]{Exercise}
\newtheorem{situation}[subsection]{Situation}
\theoremstyle{remark}
\newtheorem{remark}[subsection]{Remark}
\newtheorem{remarks}[subsection]{Remarks}
\numberwithin{equation}{subsection}
\def\lim{\mathop{\mathrm{lim}}\nolimits}
\def\colim{\mathop{\mathrm{colim}}\nolimits}
\def\Spec{\mathop{\mathrm{Spec}}}
\def\Hom{\mathop{\mathrm{Hom}}\nolimits}
\def\Ext{\mathop{\mathrm{Ext}}\nolimits}
\def\SheafHom{\mathop{\mathcal{H}\!\mathit{om}}\nolimits}
\def\SheafExt{\mathop{\mathcal{E}\!\mathit{xt}}\nolimits}
\def\Sch{\mathit{Sch}}
\def\Mor{\mathop{\mathrm{Mor}}\nolimits}
\def\Ob{\mathop{\mathrm{Ob}}\nolimits}
\def\Sh{\mathop{\mathit{Sh}}\nolimits}
\def\NL{\mathop{N\!L}\nolimits}
\def\CH{\mathop{\mathrm{CH}}\nolimits}
\def\proetale{{pro\text{-}\acute{e}tale}}
\def\etale{{\acute{e}tale}}
\def\QCoh{\mathit{QCoh}}
\def\Ker{\mathop{\mathrm{Ker}}}
\def\Im{\mathop{\mathrm{Im}}}
\def\Coker{\mathop{\mathrm{Coker}}}
\def\Coim{\mathop{\mathrm{Coim}}}
\DeclareMathSymbol{\boxtimes}{\mathbin}{AMSa}{"02}
\def\QCohstack{\mathcal{QC}\!\mathit{oh}}
\def\Cohstack{\mathcal{C}\!\mathit{oh}}
\def\Spacesstack{\mathcal{S}\!\mathit{paces}}
\def\Quotfunctor{\mathrm{Quot}}
\def\Hilbfunctor{\mathrm{Hilb}}
\def\Curvesstack{\mathcal{C}\!\mathit{urves}}
\def\Polarizedstack{\mathcal{P}\!\mathit{olarized}}
\def\Complexesstack{\mathcal{C}\!\mathit{omplexes}}
\def\Pic{\mathop{\mathrm{Pic}}\nolimits}
\def\Picardstack{\mathcal{P}\!\mathit{ic}}
\def\Picardfunctor{\mathrm{Pic}}
\def\Deformationcategory{\mathcal{D}\!\mathit{ef}}
\usepackage{cleveref}
\crefname{originalproposition}{original proposition}{original propositions}
\Crefname{originalproposition}{Original proposition}{Original propositions}
\makeatletter
\newlabel{criteria-lemma-check-representable-limit-preserving}{{07WH}{source}{}{sourcepointer-1}{}}
\newlabel{criteria-lemma-check-representable-limit-preserving@cref}{{[originalproposition][0][]07WH}{[1][0][]source}{}{}{}}
\makeatother
\crefname{originalproposition}{original proposition}{original propositions}
\Crefname{originalproposition}{Original proposition}{Original propositions}
\makeatletter
\newlabel{criteria-lemma-second-diagonal}{{07WG}{source}{}{sourcepointer-2}{}}
\newlabel{criteria-lemma-second-diagonal@cref}{{[originalproposition][0][]07WG}{[1][0][]source}{}{}{}}
\makeatother
\crefname{originalproposition}{original proposition}{original propositions}
\Crefname{originalproposition}{Original proposition}{Original propositions}
\makeatletter
\newlabel{lemma-algebraic-stack-RS}{{07WQ}{source}{}{sourcepointer-3}{}}
\newlabel{lemma-algebraic-stack-RS@cref}{{[originalproposition][0][]07WQ}{[1][0][]source}{}{}{}}
\makeatother
\crefname{originalproposition}{original proposition}{original propositions}
\Crefname{originalproposition}{Original proposition}{Original propositions}
\makeatletter
\newlabel{lemma-diagonal-representable}{{07Y4}{source}{}{sourcepointer-4}{}}
\newlabel{lemma-diagonal-representable@cref}{{[originalproposition][0][]07Y4}{[1][0][]source}{}{}{}}
\makeatother
\crefname{originalproposition}{original proposition}{original propositions}
\Crefname{originalproposition}{Original proposition}{Original propositions}
\makeatletter
\newlabel{lemma-effective}{{07X8}{source}{}{sourcepointer-5}{}}
\newlabel{lemma-effective@cref}{{[originalproposition][0][]07X8}{[1][0][]source}{}{}{}}
\makeatother
\crefname{originalproposition}{original proposition}{original propositions}
\Crefname{originalproposition}{Original proposition}{Original propositions}
\makeatletter
\newlabel{lemma-fibre-product-RS}{{07WR}{source}{}{sourcepointer-6}{}}
\newlabel{lemma-fibre-product-RS@cref}{{[originalproposition][0][]07WR}{[1][0][]source}{}{}{}}
\makeatother
\crefname{originalproposition}{original proposition}{original propositions}
\Crefname{originalproposition}{Original proposition}{Original propositions}
\makeatletter
\newlabel{lemma-fibre-product-effective}{{07X9}{source}{}{sourcepointer-7}{}}
\newlabel{lemma-fibre-product-effective@cref}{{[originalproposition][0][]07X9}{[1][0][]source}{}{}{}}
\makeatother
\crefname{originalproposition}{original proposition}{original propositions}
\Crefname{originalproposition}{Original proposition}{Original propositions}
\makeatletter
\newlabel{lemma-fibre-product-limit-preserving}{{07XM}{source}{}{sourcepointer-8}{}}
\newlabel{lemma-fibre-product-limit-preserving@cref}{{[originalproposition][0][]07XM}{[1][0][]source}{}{}{}}
\makeatother
\crefname{originalproposition}{original proposition}{original propositions}
\Crefname{originalproposition}{Original proposition}{Original propositions}
\makeatletter
\newlabel{lemma-fibre-product-tangent-spaces}{{07X2}{source}{}{sourcepointer-9}{}}
\newlabel{lemma-fibre-product-tangent-spaces@cref}{{[originalproposition][0][]07X2}{[1][0][]source}{}{}{}}
\makeatother
\crefname{originalproposition}{original proposition}{original propositions}
\Crefname{originalproposition}{Original proposition}{Original propositions}
\makeatletter
\newlabel{lemma-finite-dimension}{{07X1}{source}{}{sourcepointer-10}{}}
\newlabel{lemma-finite-dimension@cref}{{[originalproposition][0][]07X1}{[1][0][]source}{}{}{}}
\makeatother
\crefname{originalproposition}{original proposition}{original propositions}
\Crefname{originalproposition}{Original proposition}{Original propositions}
\makeatletter
\newlabel{proposition-spaces-diagonal-representable}{{07Y1}{source}{}{sourcepointer-11}{}}
\newlabel{proposition-spaces-diagonal-representable@cref}{{[originalproposition][0][]07Y1}{[1][0][]source}{}{}{}}
\makeatother
\crefname{originalproposition}{original proposition}{original propositions}
\Crefname{originalproposition}{Original proposition}{Original propositions}
\makeatletter
\newlabel{section-axioms}{{07XJ}{source}{}{sourcepointer-12}{}}
\newlabel{section-axioms@cref}{{[originalproposition][0][]07XJ}{[1][0][]source}{}{}{}}
\makeatother
\crefname{originalproposition}{original proposition}{original propositions}
\Crefname{originalproposition}{Original proposition}{Original propositions}
\makeatletter
\newlabel{stacks-lemma-2-fibre-product-gives-stack-in-setoids}{{05UI}{source}{}{sourcepointer-13}{}}
\newlabel{stacks-lemma-2-fibre-product-gives-stack-in-setoids@cref}{{[originalproposition][0][]05UI}{[1][0][]source}{}{}{}}
\makeatother
\crefname{originalproposition}{original proposition}{original propositions}
\Crefname{originalproposition}{Original proposition}{Original propositions}
\makeatletter
\newlabel{stacks-lemma-isom-as-2-fibre-product}{{04SI}{source}{}{sourcepointer-14}{}}
\newlabel{stacks-lemma-isom-as-2-fibre-product@cref}{{[originalproposition][0][]04SI}{[1][0][]source}{}{}{}}
\makeatother
\crefname{originalproposition}{original proposition}{original propositions}
\Crefname{originalproposition}{Original proposition}{Original propositions}
\makeatletter
\newlabel{stacks-lemma-stack-in-setoids-characterize}{{0432}{source}{}{sourcepointer-15}{}}
\newlabel{stacks-lemma-stack-in-setoids-characterize@cref}{{[originalproposition][0][]0432}{[1][0][]source}{}{}{}}
\makeatother
\title{The second diagonal axiom plus the axioms imply algebraicity}
\author{The Stacks Project Authors}
\date{stacks-project:artin.tex:local-snapshot-20261012; GNU FDL 1.2 or later}
\begin{document}
\maketitle
\noindent\textit{Editorial note.} Statement span is lines 2782-2796 of sources/stacks-artin/pinned-original-source.tex; proof span is lines 2798-2851, the complete author proof (52 body lines).\par\medskip
\noindent\textit{Editorial note.} Navigation stubs: 15 cited-but-not-included labels are supplied as non-load-bearing pointers with their pinned tags where the snapshot tag index has them.\par\medskip
\noindent\textit{Editorial note.} extra\_packages is empty on purpose: the retained preamble already loads xy with options and a second declaration clashes in the assembler wrapper.\par\medskip
\noindent\textit{Editorial note.} Text is retained verbatim, including the author's own typos and placeholders.\par\medskip
\noindent\textit{Editorial note.} Queue index 15 of queue\_artin.json (pb=2471, stmt 2782-2796, proof 2798-2851).\par\medskip
\noindent\textit{Editorial note.} stacks-artin wave ar01 (2026-10-14): fresh adjudication accept per lane-s3/artin\_judgments.md; pool per Chief 10:5x direct dispatch (idle-first over S12's 01:5x note, SMC precedent); artin.tex 99 rows/67 with proof judged 46 accept / 21 reject / 5 AGP-overlap avoidance rows; patch-sd-envs contract spans scanned.\par\medskip
\noindent\textit{Allowed standard prerequisites (editorial).}\begin{itemize}
\item Explicit prerequisite boundary. The cited inputs are prerequisites; the argument in the retained span is not replaced by any cited result.
\end{itemize}\medskip
\ifdefined\hypertarget\hypertarget{sourcepointer-1}{}\fi
\noindent\textit{Source reference.} Original proposition 07WH: sources/stacks-artin/pinned-original-source.tex, source lines 2851--2851. Cited by the retained statement or proof but not part of the collected spans; supplied only as a navigation pointer. The argument does not depend on its content.\par\medskip
\ifdefined\hypertarget\hypertarget{sourcepointer-2}{}\fi
\noindent\textit{Source reference.} Original proposition 07WG: sources/stacks-artin/pinned-original-source.tex, source lines 2851--2851. Cited by the retained statement or proof but not part of the collected spans; supplied only as a navigation pointer. The argument does not depend on its content.\par\medskip
\ifdefined\hypertarget\hypertarget{sourcepointer-3}{}\fi
\noindent\textit{Source reference.} Original proposition 07WQ: sources/stacks-artin/pinned-original-source.tex, source lines 743--746. Cited by the retained statement or proof but not part of the collected spans; supplied only as a navigation pointer. The argument does not depend on its content.\par\medskip
\ifdefined\hypertarget\hypertarget{sourcepointer-4}{}\fi
\noindent\textit{Source reference.} Original proposition 07Y4: sources/stacks-artin/pinned-original-source.tex, source lines 2729--2741. Cited by the retained statement or proof but not part of the collected spans; supplied only as a navigation pointer. The argument does not depend on its content.\par\medskip
\ifdefined\hypertarget\hypertarget{sourcepointer-5}{}\fi
\noindent\textit{Source reference.} Original proposition 07X8: sources/stacks-artin/pinned-original-source.tex, source lines 1323--1326. Cited by the retained statement or proof but not part of the collected spans; supplied only as a navigation pointer. The argument does not depend on its content.\par\medskip
\ifdefined\hypertarget\hypertarget{sourcepointer-6}{}\fi
\noindent\textit{Source reference.} Original proposition 07WR: sources/stacks-artin/pinned-original-source.tex, source lines 753--759. Cited by the retained statement or proof but not part of the collected spans; supplied only as a navigation pointer. The argument does not depend on its content.\par\medskip
\ifdefined\hypertarget\hypertarget{sourcepointer-7}{}\fi
\noindent\textit{Source reference.} Original proposition 07X9: sources/stacks-artin/pinned-original-source.tex, source lines 1438--1445. Cited by the retained statement or proof but not part of the collected spans; supplied only as a navigation pointer. The argument does not depend on its content.\par\medskip
\ifdefined\hypertarget\hypertarget{sourcepointer-8}{}\fi
\noindent\textit{Source reference.} Original proposition 07XM: sources/stacks-artin/pinned-original-source.tex, source lines 1693--1706. Cited by the retained statement or proof but not part of the collected spans; supplied only as a navigation pointer. The argument does not depend on its content.\par\medskip
\ifdefined\hypertarget\hypertarget{sourcepointer-9}{}\fi
\noindent\textit{Source reference.} Original proposition 07X2: sources/stacks-artin/pinned-original-source.tex, source lines 1154--1178. Cited by the retained statement or proof but not part of the collected spans; supplied only as a navigation pointer. The argument does not depend on its content.\par\medskip
\ifdefined\hypertarget\hypertarget{sourcepointer-10}{}\fi
\noindent\textit{Source reference.} Original proposition 07X1: sources/stacks-artin/pinned-original-source.tex, source lines 1030--1037. Cited by the retained statement or proof but not part of the collected spans; supplied only as a navigation pointer. The argument does not depend on its content.\par\medskip
\ifdefined\hypertarget\hypertarget{sourcepointer-11}{}\fi
\noindent\textit{Source reference.} Original proposition 07Y1: sources/stacks-artin/pinned-original-source.tex, source lines 2641--2649. Cited by the retained statement or proof but not part of the collected spans; supplied only as a navigation pointer. The argument does not depend on its content.\par\medskip
\ifdefined\hypertarget\hypertarget{sourcepointer-12}{}\fi
\noindent\textit{Source reference.} Original proposition 07XJ: sources/stacks-artin/pinned-original-source.tex, source lines 2462--2489. Cited by the retained statement or proof but not part of the collected spans; supplied only as a navigation pointer. The argument does not depend on its content.\par\medskip
\ifdefined\hypertarget\hypertarget{sourcepointer-13}{}\fi
\noindent\textit{Source reference.} Original proposition 05UI: sources/stacks-artin/pinned-original-source.tex, source lines 2851--2851. Cited by the retained statement or proof but not part of the collected spans; supplied only as a navigation pointer. The argument does not depend on its content.\par\medskip
\ifdefined\hypertarget\hypertarget{sourcepointer-14}{}\fi
\noindent\textit{Source reference.} Original proposition 04SI: sources/stacks-artin/pinned-original-source.tex, source lines 2851--2851. Cited by the retained statement or proof but not part of the collected spans; supplied only as a navigation pointer. The argument does not depend on its content.\par\medskip
\ifdefined\hypertarget\hypertarget{sourcepointer-15}{}\fi
\noindent\textit{Source reference.} Original proposition 0432: sources/stacks-artin/pinned-original-source.tex, source lines 2851--2851. Cited by the retained statement or proof but not part of the collected spans; supplied only as a navigation pointer. The argument does not depend on its content.\par\medskip
\begin{proposition}
\label{proposition-second-diagonal-representable}
Let $S$ be a locally Noetherian scheme. Let
$p : \mathcal{X} \to (\Sch/S)_{fppf}$ be a category fibred in groupoids.
Assume that
\begin{enumerate}
\item $\Delta_\Delta : \mathcal{X} \to
\mathcal{X} \times_{\mathcal{X} \times \mathcal{X}} \mathcal{X}$
is representable by algebraic spaces,
\item $\mathcal{X}$ satisfies axioms [-1], [0], [1], [2], [3], [4], and [5]
(see Section \ref{section-axioms}),
\item $\mathcal{O}_{S, s}$ is a G-ring for all finite type points $s$ of $S$.
\end{enumerate}
Then $\mathcal{X}$ is an algebraic stack.
\end{proposition}

\begin{proof}
We first prove that $\Delta : \mathcal{X} \to \mathcal{X} \times \mathcal{X}$
is representable by algebraic spaces. To do this it suffices to show
that
$$
\mathcal{Y} =
\mathcal{X} \times_{\Delta, \mathcal{X} \times \mathcal{X}, y} (\Sch/V)_{fppf}
$$
is representable by an algebraic space for any affine scheme $V$ locally
of finite presentation over $S$ and object $y$ of
$\mathcal{X} \times \mathcal{X}$ over $V$, see
Criteria for Representability, Lemma
\ref{criteria-lemma-check-representable-limit-preserving}\footnote{The
set theoretic condition in Criteria for Representability, Lemma
\ref{criteria-lemma-check-representable-limit-preserving}
will hold: the size of the algebraic space $Y$ representing $\mathcal{Y}$ is
suitably bounded. Namely, $Y \to S$ will be locally of finite type and $Y$
will satisfy axiom [-1]. Details omitted.}.
Observe that $\mathcal{Y}$ is fibred in setoids
(Stacks, Lemma \ref{stacks-lemma-isom-as-2-fibre-product})
and let $Y : (\Sch/S)_{fppf}^{opp} \to \textit{Sets}$,
$T \mapsto \Ob(\mathcal{Y}_T)/\cong$ be the functor of isomorphism
classes. We will apply
Proposition \ref{proposition-spaces-diagonal-representable}
to see that $Y$ is an algebraic space.

\medskip\noindent
Note that
$\Delta_\mathcal{Y} : \mathcal{Y} \to \mathcal{Y} \times \mathcal{Y}$
(and hence also $Y \to Y \times Y$)
is representable by algebraic spaces by condition (1) and
Criteria for Representability, Lemma \ref{criteria-lemma-second-diagonal}.
Observe that $Y$ is a sheaf for the \'etale topology by
Stacks, Lemmas \ref{stacks-lemma-stack-in-setoids-characterize} and
\ref{stacks-lemma-2-fibre-product-gives-stack-in-setoids}, i.e.,
axiom [0] holds. Also $Y$ is limit preserving by
Lemma \ref{lemma-fibre-product-limit-preserving}, i.e., we have [1].
Note that $Y$ has (RS), i.e., axiom [2] holds, by
Lemmas \ref{lemma-algebraic-stack-RS} and
\ref{lemma-fibre-product-RS}. Axiom [3] for $Y$ follows
from Lemmas \ref{lemma-finite-dimension} and
\ref{lemma-fibre-product-tangent-spaces}.
Axiom [4] follows from Lemmas \ref{lemma-effective} and
\ref{lemma-fibre-product-effective}.
Axiom [5] for $Y$ follows directly from openness of versality
for $\Delta_\mathcal{X}$ which is part of axiom [5] for $\mathcal{X}$.
Thus all the assumptions of
Proposition \ref{proposition-spaces-diagonal-representable}
are satisfied and $Y$ is an algebraic space.

\medskip\noindent
At this point it follows from Lemma \ref{lemma-diagonal-representable}
that $\mathcal{X}$ is an algebraic stack.
\end{proof}

\end{document}
